\section{Related Work}
\label{sec:related_work}

We situate our work at the intersection of static model merging, test-time/dynamic merging, and a critical analysis of quantum-inspired deep learning.

\subsection{Static Model Merging}
Model merging operates in weight space to fuse specialized, independently fine-tuned task experts back into a cohesive multi-task model without additional heavy training~\cite{wortsman2022model}. Early methods proposed simple uniform parameter averaging or task arithmetic (linear addition of task-specific parameter offsets from a pre-trained base)~\cite{ilharco2022editing}. Although task arithmetic works well for highly related tasks, it suffers from parameter interference. To resolve coordinate-wise disagreements, Yadav et al. introduced TIES-Merging~\cite{yadav2023ties}, which applies sign consensus voting and magnitude pruning. Sparse Task Arithmetic (STA)~\cite{sta} later demonstrated that sign-agreement heuristics are often redundant, and simple uniform layer-wise magnitude pruning of task vectors is sufficient. To calibrate these merging coefficients, offline validation frameworks like OFS-Tune~\cite{ofstune} and SuiteMerge~\cite{suitemerge} utilize tiny, labeled validation sets with low-degree polynomial constraints to optimize static layer-wise weights, demonstrating that zero test-time compute can achieve high multi-task performance under stable distributions.

\subsection{Test-Time and Dynamic Model Merging}
In dynamic, non-stationary deployment environments, static coefficients are sub-optimal. Test-Time Adaptation (TTA) of merging coefficients addresses this by adapting weights on unlabeled test streams. AdaMerging~\cite{adamerging} dynamically optimizes merging coefficients at test-time by minimizing the entropy of the merged model's predictions over the incoming data stream. ZipMerge~\cite{zipmerge} extends this by jointly optimizing merging coefficients and parameter pruning masks on edge devices. PolyMerge~\cite{polymerge} proposes fitting continuous polynomial trajectories to task vectors to adapt the merging weights continuously. However, these methods assume test streams are homogeneous and stable. In real-world deployment, unlabeled streams can be highly heterogeneous, leading to high variance and optimization instabilities. 

\subsection{Quantum-Inspired Deep Learning and the Baseline Confounder}
In recent years, the machine learning community has witnessed a surge in ``quantum-inspired'' deep learning architectures. These works often wrap standard neural operations, such as attention or gating, in the language of quantum mechanics (e.g., modeling inputs as wavefunctions, states as density matrices, or projections as operators/collapses)~\cite{qwsmerge}. While these analogies are mathematically elegant, they frequently suffer from what we term the \textit{baseline confounder}: they compare their complex, highly-tuned quantum formulations against crippled, unregularized classical baselines, attributing the performance improvements to ``quantum-like wave interference'' rather than basic structural design. 

For example, the recent Quantum Wavefunction Superposition Merging (QWS-Merge) SOTA~\cite{qwsmerge} introduces an input-dependent layer-wise cosine wave routing network. QWS-Merge is compared against a ``Linear Router'' baseline that is intentionally restricted to be global (instead of layer-wise) and is evaluated under a low-data calibration regime with zero regularization (no weight decay). In this paper, we break this confounding loop by setting up a mathematically identical, low-dimensional classical comparison that is properly regularized, exposing the redundancy of the quantum wave-interference analogy.
